\chapter*{Заключение} \label{ch-conclusion}
\addcontentsline{toc}{chapter}{Заключение}	% в оглавление 

В ходе выполнения лабораторной работы было разработано консольное приложение на языке Java для фильтрации IP-адресов по чёрному списку (blacklist) с использованием двух алгоритмов (детерминированный HashSet и вероятностный Bloom Filter), а также реализован и протестирован набор модульных тестов с последующей проверкой их качества методом мутационного тестирования.

В работе спроектировано и реализовано консольное приложение с меню, поддерживающее два режима работы: интерактивную проверку IP-адресов и пакетную обработку данных из JSON-файла с формированием JSON-отчёта. Архитектура приложения построена модульно: отдельно выделены пользовательский интерфейс (Main, ConsoleApp), модуль фильтрации (com.ipfilter.filter) и модуль сервисов (com.ipfilter.service). Для решения задачи реализованы два конкурирующих алгоритма фильтрации: детерминированный подход на основе HashSet (100\% точность, высокое потребление памяти) и вероятностный подход на основе Bloom Filter (возможны False Positive, значительная экономия памяти), что позволило использовать их как для сравнения производительности, так и для формирования рекомендаций пользователю.

Для тестирования были применены четыре техники проектирования тестов: две техники, основанные на спецификации (Boundary Value Analysis, Equivalence Partitioning), и две техники, основанные на структуре программы (Branch Testing, Data Flow Testing). На их основе были сформированы контрольные примеры для позитивных и негативных сценариев, проверяющие входные данные, выходные результаты и ключевые внутренние условия обработки. В рамках тестирования были покрыты проверки обоих алгоритмов фильтрации, сервис обработки данных, пакетная обработка JSON-данных и запись отчёта. Также реализовано сравнение результатов работы двух алгоритмов, что повысило надёжность тестирования основной логики.

В качестве основной библиотеки модульного тестирования использован стек JUnit 5 (вместе с Hamcrest и Mockito). Его основными преимуществами в рамках данной работы стали: удобная структура тестов, поддержка параметризованных тестов, богатый набор assertions и assumptions, интеграция с IDE и Maven, а также возможность использования мокирования для изоляции зависимостей. Дополнительно Hamcrest упростил запись читаемых проверок, а Mockito позволил протестировать взаимодействия и зависимости без усложнения основной логики. К недостаткам можно отнести более высокий порог входа по сравнению с базовыми тестами (особенно при использовании моков и параметризации), а также необходимость дополнительной настройки для мутационного тестирования (Pitest + JUnit5 plugin).

В результате работы был сформирован тестовый набор, покрывающий как позитивные, так и негативные сценарии для основных классов приложения. Качество тестов было дополнительно проверено с помощью Pitest. По итоговым результатам мутационного тестирования достигнуты следующие показатели:
\begin{itemize}
	\item Line Coverage - 78\%
	\item Mutation Coverage - 80\%
	\item Test Strength - 93\%
\end{itemize}
Полученные значения метрик свидетельствуют о приемлимом качестве тестового набора и позволяют считать его практически достаточным для учебной задачи.

Таким образом, поставленные в работе цели были достигнуты: приложение спроектировано и реализовано, основная функциональность покрыта модульными тестами, тесты спроектированы с использованием четырёх техник, а их качество подтверждено метриками и результатами мутационного тестирования. Использование мутационного тестирования позволило не только получить количественные метрики качества, но и выявить конкретные слабые места в тестовом наборе для последующего улучшения.
